\documentclass[11pt]{article}
\usepackage[margin=1in]{geometry}
\usepackage{amsmath,amssymb,amsthm}
\newtheorem{theorem}{Theorem}
\newtheorem{lemma}{Lemma}
\begin{document}
\title{Laboratory V84: Exact $\mathrm{FP}^{\mathrm{NP}}$ Girth Extraction and a Hall-Expander Promise Reduction}
\author{NC$^0_k$-Avoid Laboratory}
\date{August 2026}
\maketitle

Let $B=(M,X;E)$ present a transversal matroid, with $|M|=m$ and
$|N(e)|\leq 3$ for every $e\in M$. Let $g(B)$ be its girth and let the oracle
language ask whether $g(B)\leq L$.

\begin{theorem}[Exact extraction]
When dependence is guaranteed, $g(B)$ is computable with
$\lceil\log_2 m\rceil$ oracle queries. A fixed-order deletion self-reduction
then returns a canonical shortest circuit using at most $m$ additional oracle
queries. Without guaranteed dependence, one existence query suffices.
Every query preserves maximum left degree three.
\end{theorem}

\begin{proof}
Binary-search the least accepted threshold. For the self-reduction, process
left elements in a fixed order and delete $e$ whenever the oracle accepts the
restriction $B\setminus e$ at the already-known threshold $g(B)$. Deletion
cannot create a shorter circuit, so every accepted deletion preserves girth.
At termination every remaining element belongs to every shortest circuit;
since a shortest circuit is contained in the remaining ground set, the
remaining set is that unique circuit.
\end{proof}

\begin{lemma}[Exact Hall witness]
If $C$ is a transversal circuit, then $|N(C)|=|C|-1$. Hence a shortest circuit
of size $g$ yields the exact minimum Hall-neighborhood value $h^*=g-1$.
\end{lemma}

\begin{proof}
Circuit dependence gives $|N(C)|<|C|$, because any proper Hall-deficient subset
would contradict minimality. For each $e\in C$, the independent set
$C\setminus\{e\}$ has at least $|C|-1$ neighbors, all in $N(C)$. Thus equality
holds. Every deficient set contains a circuit, so no smaller neighborhood
than $g-1$ is possible.
\end{proof}

\begin{theorem}[Local avoidance or Hall expansion]
For every threshold $L$, an $\mathrm{FP}^{\mathrm{NP}}$ preprocessing procedure
returns either:
\begin{enumerate}
\item an avoided output, using $2^{g-1}\operatorname{poly}(n+m)$ deterministic
local work when $g\leq L$; or
\item the promise $|N(S)|\geq |S|$ for every $S\subseteq M$ with $|S|\leq L$
when $g>L$.
\end{enumerate}
For fixed $c$ and $L=c\log_2(n+m)$, the first branch is polynomial.
\end{theorem}

\begin{proof}
For a shortest circuit $C$, the local map on $N(C)$ has at most
$2^{|N(C)|}=2^{g-1}$ images inside a codomain of size $2^g$. Enumerate the
local range, choose a missing projection, and lift it by fixing arbitrary bits
outside $C$. If $g>L$, every set of at most $L$ outputs is independent, and
Hall's theorem gives the stated noncontraction promise.
\end{proof}

This is a promise reduction, not a complete solver for the large-girth branch.
It does not establish matrix rigidity, circuit lower bounds, or $P\neq NP$.
\end{document}
